\section{问题四模型的建立与求解}

\subsection{问题四的描述与分析}

前三问依次释放计算柔性和能源柔性：问题二允许任务迁移与延迟但固定能源动作，问题三固定计算负荷只优化储能；问题四要求\textbf{同时优化任务与能源--储能系统}。附件中 RT 的服务边界严格，而 AI Training 贡献约 80.14\% GPU-hour 和 87.12\% AI IT 能量；区域的网络、算力、PUE、储能和售电优势又彼此错位。更关键的是，迁入任务与储能充电会竞争同一新能源余量，因此问题二的固定能源边际并不等于储能重优化后的真实边际。

因此将任务放置直接映射为 AI IT 与设施负荷，并对每个候选负荷状态重新优化能源和储能。约 $2.33\times10^8$ 个合法任务放置由区域多割 Benders 与完整域列生成联合处理；服务质量则在成本最优面内依次压缩总等待和网络时延。工作负载转移与能源系统需要联合管理的思想与近期数据中心 integrated workload--energy management 研究一致\cite{guo_integrated_2025}。总体流程见图\ref{fig:q4-flow}。

\begin{figure}[H]
\centering
\includegraphics[width=\textwidth]{\ProjectRoot/modules/50_q4/figures/q4_flow.png}
\caption{问题四计算--能源联合优化总体流程}
\label{fig:q4-flow}
\end{figure}

题面要求同时考察 Cost、Carbon、Latency、QoS、RenewableUtilization 与 regional PeakNetImport。本文不把六项指标等权压缩为单一分数：基准场景以 Cost 为第一经济目标，在其最优面内依次压缩总等待和网络时延；Carbon 在碳预算场景中作为 $\varepsilon$-constraint 激活，新能源利用率与区域峰值净购电则保持题面统一口径作为系统评价量。这样避免在指标已退化或高度同向时人为制造权重冲突。

\subsection{预备工作：联合优化必要性与顺序失真}

50000 个任务数量近似三等分，但资源权重和柔性差异明显：Training、Batch、RT 分别约占 80.14\%、15.32\%、4.54\% GPU-hour；RT 到达即开工且合法区域最少。RegionD 的 AvailableGPU 最大（1472），D/E/F 的 PUE 更低且储能、售电能力更强，而 A/B/C 网络时延更有利。六区域逐时 $\mathrm{AvailableRenewable}$ 又完全相同，所以区域绿色优势不能由标签预设，只能由 PUE、负荷、价格、碳强度、储能和售电边界共同决定。

给定任务排程 $\mathbf x$，动态新能源盈余为
\begin{equation}
H_{rt}(\mathbf x)=\mathrm{AvailableRenewable}_{rt}-L_{rt}^{Fac}(\mathbf x).
\end{equation}
同一 1 MWh 新能源可被迁入任务当期消纳，也可先充入储能跨时释放，两类柔性既互补又竞争。

顺序优化探针进一步验证了这一点：选取 3 个仅能在 E/F 执行的真实 RT 任务，共 8 种放置。固定能源动作时，“全部迁 F”可节约约 3613.67 CNY；但分别重求完整能源--储能 LP 后，8 种放置的最优能源成本相同，此时保留在 E 的方案网络时延更低。因此
\begin{equation}
\text{固定能源边际价值}\neq\text{储能重优化后的真实边际价值}.
\label{eq:q4-sequential-failure}
\end{equation}
能源 LP 的负荷平衡对偶与真实边际成本排序仍高度相关（Pearson 约 0.985、Spearman 约 0.898），适合形成累积 Benders 支撑超平面，而不适合只保留单次线性近似。

\subsection{IWESS 计算--能源--储能联合调度模型}

本文将上述联合模型记为 \textbf{Integrated Workload--Energy--Storage Scheduling（IWESS）}。该名称用于概括任务时空决策与 BESS/电网连续再优化构成的联合数学对象，不把缩写本身作为创新声明。对任务 $i$ 按 SLA、Arrival、LatestFinish 与 2406 边界构造完整合法放置集 $\Omega_i$，仅对 $(r,s)\in\Omega_i$ 定义 $x_{irs}\in\{0,1\}$。任务与小时的重叠量为
\begin{equation}
\omega_{ist}=\max\bigl(0,\min(t+1,s+p_i)-\max(t,s)\bigr),
\end{equation}
从而
\begin{align}
L^{AI}_{rt}(\mathbf x)&=\sum_{i,s}\alpha_{k_i}g_i\omega_{ist}x_{irs},\label{eq:q4-ai-load}\\
L^{Fac}_{rt}(\mathbf x)&=\mathrm{PUE}_r\bigl(L^{NonAI}_{rt}+L^{AI}_{rt}(\mathbf x)\bigr).\label{eq:q4-fac-load}
\end{align}

能源侧沿用问题三 BESS-ED 变量 $\mathbf y_r=(u,qR,qG,d,gL,s,w,E)$。对给定设施负荷 $L_r^{Fac}$，区域可行集 $\mathcal Y_r(L_r^{Fac})$ 由新能源平衡、负荷平衡、SOC 递推、充放电功率、购售电和终端 SOC 约束共同定义。联合目标为
\begin{align}
\mathrm{Cost}(\mathbf y)&=\sum_{r,t}\bigl(c^{buy}_{rt}G^{buy}_{rt}-c^{sell}_{rt}s_{rt}\bigr),\label{eq:q4-cost}\\
W(\mathbf x)&=\sum_{i,r,s}(s-a_i)x_{irs},\qquad
D(\mathbf x)=\sum_{i,r,s}\ell_{o_i r}x_{irs}.\label{eq:q4-qos}
\end{align}

\subsubsection{联合模型汇总}

IWESS 的基准字典序目标直接展开为
\begin{equation}
\operatorname*{lex\,min}_{\mathbf x,\mathbf y}\quad
\left(
\sum_{r,t}\bigl(c^{buy}_{rt}G^{buy}_{rt}-c^{sell}_{rt}s_{rt}\bigr),
\quad
\sum_{i,r,s}(s-a_i)x_{irs},
\quad
\sum_{i,r,s}\ell_{o_i r}x_{irs}
\right)
\label{eq:q4-summary-obj}
\end{equation}
\begin{equation}
\mathrm{s.t.}\quad
\left\{
\begin{aligned}
&\sum_{(r,s)\in\Omega_i}x_{irs}=1,&&\forall i,\\
&\sum_{i,s}g_i\omega_{ist}x_{irs}\le\mathrm{AvailableGPU}_r,&&\forall r,t,\\
&L^{IT}_{rt}(\mathbf x)\le\mathrm{MaxITPower}_r,\quad
L^{Fac}_{rt}(\mathbf x)\le\mathrm{MaxFacilityPower}_r,&&\forall r,t,\\
&\mathbf y_r\in\mathcal Y_r\bigl(L_r^{Fac}(\mathbf x)\bigr),&&\forall r,\\
&x_{irs}\in\{0,1\},&&(r,s)\in\Omega_i.
\end{aligned}
\right.
\label{eq:q4-summary-con}
\end{equation}
其中 $\Omega_i$ 已编码网络 SLA、Arrival、LatestFinish、RT 即时开工和 finish$\le2406$；$\mathcal Y_r$ 内含问题三全部能源、储能与电网物理约束。

\subsection{模型求解：Benders--CG 与字典序精化}

给定任务负荷后，区域最小能源成本为
\begin{equation}
Q_r(L_r)=\min_{\mathbf y_r\in\mathcal Y_r(L_r)}\sum_t\left(c^{buy}_{rt}G^{buy}_{rt}-c^{sell}_{rt}s_{rt}\right),
\end{equation}
基准成本场景无全局碳预算，故 $Q(L)=\sum_rQ_r(L_r)$。设第 $k$ 个负荷点的负荷平衡对偶为 $\lambda_{rt}^k$，区域最优性割为
\begin{equation}
\theta_r\ge Q_r(L_r^k)+\sum_t\lambda_{rt}^k\left(L_{rt}^{Fac}(\mathbf x)-L_{rt}^{Fac}(\mathbf x^k)\right).
\label{eq:q4-benders-cut}
\end{equation}

完整合法域约含 $2.33\times10^8$ 个放置，故只保留活动列并对完整隐式域做精确定价，不使用 Top-k 或人工等待窗。Benders decomposition 与 column generation 的组合是处理大规模离散调度与分解结构的成熟求解路线，近期 Omega 研究亦给出了二者联合用于大规模 allocation/scheduling 的实例\cite{wang_bcg_2026}。附件满足 $\mathrm{MaxFacility}=\mathrm{PUE}\times\mathrm{MaxIT}$ 且 $\mathrm{FacilityLoad}=\mathrm{PUE}\times\mathrm{ITLoad}$，因此实现中无损删除重复 Facility 行；Benders 割按非零对偶稀疏存储。40-task 显式校准含 3139 个合法放置，列生成仅用 44 个活动列即可得到相同根节点目标值。

字典序求解分三层进行。第一层最小化 Cost；第二层在 $\mathrm{Cost}\le C^*+\varepsilon_C$ 内最小化总等待；第三层固定前两层锚点后最小化总网络时延。每一层均重新执行完整域定价和真实能源复核；若后序阶段产生新列，则在最终扩展列池上从第一层重新认证。算法见算法\ref{alg:q4-benders-cg}。

% 算法: Benders 分解 + 精确列生成 + 六指标锚点折中
\begin{algorithm}[H]
\setstretch{0.92}
\SetAlFnt{\normalsize}
\SetArgSty{textnormal}
\SetKwInput{KwIn}{Input}
\SetKwInput{KwInit}{Initialize}
\SetKwInput{KwOut}{Output}
\SetKwFor{For}{for}{do}{end for}
\SetKwFor{ForEach}{for each}{do}{end for}
\SetKwFor{While}{while}{do}{end while}
\SetKwIF{If}{ElseIf}{Else}{if}{then}{else if}{else}{end if}
\SetKw{KwRet}{return}
\caption{Benders 分解、精确列生成与六指标折中求解}
\label{alg:q4-benders-cg}
\KwIn{任务集 $\mathcal{I}$，区域 GPU/IT 容量，新能源出力，BESS 参数，购售电价，SLA，时域及场景参数}
\KwInit{构建可行参考排程 $\mathbf{x}^0$ 与初始列池 $\mathcal{C}$；初始化 Benders 割集}
\BlankLine
\emph{阶段一：基准联合模型的 Benders--CG 闭合}\;
\While{仍存在违反的 Benders 割或负 reduced-cost 列}{
    在当前列池与割集上求 Restricted Master LP，得到 $\mathbf{x}^*$ 与 Master dual\;
    \ForEach{区域 $r$}{
        由 $\mathbf{x}^*$ 重算 $L_{rt}^{\mathrm{Fac}}$，求解真实能源/储能 LP\;
        若 recourse 值违反当前 $\theta_r$，加入区域 Benders 最优性割\;
    }
    若本轮无新割，则对每个任务在其\textbf{完整合法域}上执行 exact pricing\;
    对 Batch/Training 使用 $(k_i,r,p_i)$ 模板与 suffix-min；RT 固定到达时刻、枚举全部时延可行区域\;
    将所有满足定价准则的负 reduced-cost 列加入 $\mathcal{C}$\;
}
\BlankLine
\emph{阶段二：整数精化与真实 recourse 复核}\;
在闭合列池上求 restricted integer master，得到可行整数解与上界 $UB$\;
以 root LP 给出下界 $LB$，计算 integrality gap\;
对整数解重新调用真实 Energy LP，并复核 GPU/IT、SLA、LatestFinish、2406、SOC 与购售电边界\;
\If{integrality gap 不满足门限}{
    仅在必要时进入 branch-and-price / branch-price-and-Benders；不得以经验 top-$k$ 删列代替 exact search\;
}
\BlankLine
\emph{阶段三：六个单目标锚点}\;
\ForEach{$j\in\{C,E,L,Q,R,P\}$}{
    在同一联合模型和完整硬约束下求单目标锚点 $\mathbf{x}^{(j)}$\;
    记录六指标 payoff 向量 $\mathbf{f}(\mathbf{x}^{(j)})$\;
}
由 payoff table 构造 $f_j^{I}$ 与 $f_j^{N}$，定义无量纲退让
$d_j=(f_j-f_j^I)/(f_j^N-f_j^I)$\;
\BlankLine
\emph{阶段四：标准化最坏遗憾折中}\;
求解 $\min z$，满足 $d_j\le z$ 及全部联合调度硬约束，得到 $z^*$\;
在 $z\le z^*+\varepsilon$ 内最小化 $\sum_j d_j$\;
若仍有数值等价解，仅作解释性 tie-break：先减小最大绝对等待，再减少迁移任务数\;
\BlankLine
\emph{阶段五：场景重优化}\;
\ForEach{碳预算、电价机制或新能源波动场景 $s$}{
    复用基准 active column pool 与整数 incumbent 作为 warm start\;
    按场景参数重建必要的 Energy LP 与 Benders cuts，并重新执行 exact pricing\;
    使用与基准场景完全相同的六指标折中规则重新优化；\textbf{不得固定基准/QoS 排程只重算能源}\;
    记录 Cost、Carbon、Latency、QoS、RenewableUtilization、PeakNetImport 与迁移/BESS 策略变化\;
}
\KwRet{基准六指标折中整数解、六个锚点、各场景重优化结果及完整审计}
\KwOut{任务区域--时间分配、BESS 充放电、购售电、新能源利用、六指标与场景比较}
\end{algorithm}


\subsection{结果与分析}

40-task 精确校准中，仅最小成本可出现 38/40 个任务迁移、总等待 172 h；加入等待与时延目标后得到迁移 5 个、总等待 6 h、最大等待 2 h、总时延 437 ms 的精确字典序方案，说明服务质量精化确有必要。

50k 最终扩展列池包含 137158 个活动列和 1173 条 Benders 割，两轮再认证结果完全一致。

\begin{table}[!htbp]
\centering
\caption{50k 任务最终扩展列池的字典序再认证结果}
\label{tab:q4-final-recert}
\begin{tabular}{lccc}
\toprule
\rowcolor{CumcmLightBlue}
指标 & 第1轮 & 第2轮 & 最终代表值 \\
\midrule
真实能源成本/CNY & $-459{,}340{,}688.8007$ & $-459{,}340{,}688.8007$ & $-459{,}340{,}688.8007$ \\
总等待/h & 0 & 0 & 0 \\
总网络时延/ms & 250,860 & 250,860 & 250,860 \\
活动列数 & 137,158 & 137,158 & 137,158 \\
Benders 割数 & 1,160 & 1,173 & 1,173 \\
最小缺失列约化成本 & 0 & 0 & 0 \\
最大区域 Benders 违反/CNY & 0 & 0 & 0 \\
\bottomrule
\end{tabular}
\end{table}

最终排程覆盖全部 50000 个 TaskID，等待均为 0；仅 61 个任务跨区，迁移率 0.122\%，其中 Training 52、Batch 8、RT 1。总网络时延 250860 ms，平均 5.0172 ms，P95 为 5 ms，最大 58 ms，最晚完成 2405.6 h；任务、容量、SLA、截止时间、能源平衡、SOC 和购售电边界全部通过审计，最大物理残差约 $1.14\times10^{-13}$。

\subsubsection{储能柔性替代计算柔性}

问题二固定能源动作时需迁移 74.896\% 任务，平均等待 20.475 h、最大等待 2016 h；问题四联合优化后总等待降为 0，迁移率仅 0.122\%。这说明在本附件数据下，\textbf{储能时间柔性替代了绝大部分原本由任务长时间平移和大规模迁移承担的能源调节}，计算侧只需少量空间修正即可维持成本最优面内的服务质量。

问题三与问题四当前绝对成本差约 $1.23\times10^5$ CNY，但第 2406 h 跨问结算口径尚待统一，因此暂不把该差额解释为正式联合增量收益。

\subsection{模型检验与正式场景}

最终两轮均无缺失负约化成本列，区域 Benders 违反为 0，最终活动列池 MIP gap 为 0；但这不能自动推出完整域整数全局最优。时延层完整域 LP 下界为 250469.9546 ms，最终整数值 250860 ms，相对差
\begin{equation}
\frac{250860-250469.9546}{250860}\approx0.1555\%.
\end{equation}
因此正式表述为：\textbf{完整域 LP 定价与区域 Benders 已闭合，并获得真实能源可行、两轮再认证稳定的最优已知字典序整数方案；不声称完整域整数全局最优。}

不同碳约束、电价机制和新能源波动场景均沿用同一 IWESS 数学模型，只改变外生参数或约束并重新认证相应策略；前期 $\gamma=1.4$ 新能源压力实验只用于机制诊断，不作为正式场景结论。